\begin{center}
\begin{minipage}{\linewidth}
\centering
\textbf{垂直切线的交点在同一个圆上}\par\smallskip
\begin{tikzpicture}[scale=1.1, every node/.style={font=\small}, line cap=round]
\draw[->,gray] (-4,0)--(4.2,0) node[right] {$x$};
\draw[->,gray] (0,-3.8)--(0,4) node[above] {$y$};
\draw[CStatement,thick] (0,0) ellipse [x radius=3,y radius=2];
\node[below left] at (0,0) {$O$};
\draw[CExplanation,dashed,thick] (0,0) circle [radius={sqrt(13)}];

\draw[CProof,thick] (3,-.65)--(3,3.2) node[above] {$x=a$};
\draw[CProof,thick] (-.7,2)--(3.65,2) node[right] {$y=b$};
\draw[CExample,thick] (0,0)--(3,2);
\draw (2.76,2)--(2.76,1.76)--(3,1.76);
\fill (3,2) circle(2pt) node[above right] {$P(a,b)$};
\node[below right] at (3,0) {$a$}; \node[above left] at (0,2) {$b$};
\node[fill=white,inner sep=3pt] at (0,-4.2) {蒙日圆：$x^2+y^2=a^2+b^2$};

\end{tikzpicture}
\par 椭圆实线、轨迹圆虚线；图取 $a=3,b=2$，全图等比例。
\end{minipage}
\end{center}
